% !TeX root = ../数值分析与科学计算笔记.tex
% 由主文档通过 \input 载入；请编译主文档。
\setcounter{chapter}{0} % 保留原书章号。
\chapter{绪论：从数学问题到可靠计算}
\label{chap:na-intro}

\section{研究对象与基本问题}
\emph{数值分析（numerical analysis）}研究数学问题的近似计算方法及其误差、稳定性和计算代价。\emph{科学计算（scientific computing）}进一步把这些方法组织成能够解决科学与工程问题的计算过程，涉及模型、离散化、算法实现和结果验证。一个方程有唯一解，不等于有限精度计算能够准确地求出它；一个程序返回许多位小数，也不等于这些小数具有相应的有效精度。

以求解导热问题为例，物理规律首先给出温度满足的微分方程，边界条件和材料参数决定具体模型；空间网格把连续函数转化为有限维向量，时间步进把连续演化转化为递推，线性代数算法再求出每一步的近似值。测量误差、模型误差、离散误差、迭代误差和舍入误差发生在不同层次，应分别识别。减小网格宽度通常降低离散误差，却也增加计算代价，并可能使代数系统更难求解。

本书围绕三个相互联系的问题展开：所求数学对象是否存在、唯一并且对数据连续依赖；离散或迭代近似是否趋于该对象以及趋近得多快；有限精度实现是否保留了理论算法的可靠性。第一类问题属于问题本身的性质，后两类还取决于算法。选择算法时应同时考察精度要求、数据规模和可利用的结构，而不能只比较公式的阶数。

\begin{example}[很小的残差能说明什么]
\label{ex:na-intro-residual}
设 $0<\varepsilon\ll1$，考虑 $\varepsilon x=\varepsilon$，其精确解为 $x=1$。近似解 $\widehat x=0$ 的残差为 $\varepsilon$，绝对值很小，但解的绝对误差等于 $1$。这说明绝对残差必须结合问题尺度解释。若把等式两端同除以 $\varepsilon$，残差变成 $1$，而数学问题没有改变。后文将用归一化残差和条件数把这一直觉写成可核查的不等式。
\end{example}

\section{本科定位、知识依赖与阅读路线}
本书面向已学习一元与多元微积分、线性代数的本科生。常微分方程初值问题的基本概念在相应章节回顾；蒙特卡洛方法还需要期望、方差、独立性和中心极限定理。编程部分仅要求理解变量、循环、函数和数组。线性空间的完备性、矩阵谱分解等结果按需要说明，不要求预先学习泛函分析。

第~\ref{chap:na-errors} 章建立误差语言，第~\ref{chap:na-roots} 章讨论非线性方程；第~\ref{chap:na-direct} 至第~\ref{chap:na-eigen} 章组成数值线性代数单元；第~\ref{chap:na-interpolation}、\ref{chap:na-quadrature} 章研究函数的离散表示与数值微积分；第~\ref{chap:na-ode}、\ref{chap:na-pde} 章把这些方法用于微分方程；第~\ref{chap:na-fourier}、\ref{chap:na-montecarlo} 章介绍傅里叶计算和随机积分；第~\ref{chap:na-practice} 章组织计算实验。最小二乘中的函数逼近依赖正交投影，插值则依赖节点值约束，两者采用不同的误差准则。

一学期课程可以把误差、求根、直接法、最小二乘、基本迭代、插值、数值积分和常微分方程作为主干，按专业需要选讲特征值、偏微分方程、傅里叶和蒙特卡洛专题。完整阅读适合配合持续的上机练习。正文中的例题展示方法与推导，章末习题包括证明、手算、反例和实验；附录给出解答或实验判据，实验中的预期现象不冒充实际运行结果。

\section{编写依据与资料使用范围}
课程范围参考三个可公开核验的来源。清华大学出版社对喻文健《数值分析与算法》第三版的介绍覆盖误差、求根、线性方程组、特征值、逼近、积分和常微分方程，并强调算法与应用\cite{yu2020}；北京大学公布的本科数值分析大纲包含函数逼近、数值微积分、非线性方程、常微分方程、快速傅里叶变换和蒙特卡洛方法\cite{pku-course}；MIT 18.330 的 2012 年公开课程突出离散化、收敛速度、傅里叶分析与数值实验，并说明其线性代数内容较少\cite{mit-syllabus,mit-notes}。

本书据此采用覆盖基础数值方法的主线，补充两类科学计算专题。上述资料用于课程定位、主题核对和延伸阅读，正文、例题与习题独立组织；本书并非上述学校的官方教材、原教材译本或逐章改写。由于只核验公开大纲与图书介绍，不把未查阅的纸质教材页码或具体证明列为引文依据。书中较深结果若不展开完整证明，会明确标示证明范围。

\section{记号与算法约定}
实数和复数分别记为 $\mathbb R$、$\mathbb C$，$\mathrm i^2=-1$。向量通常为列向量，$A^{\mathsf T}$ 表示转置，$A^*$ 表示共轭转置，$I$ 为适当阶数的单位矩阵。$\mathcal P_n$ 表示次数不超过 $n$ 的实多项式空间。$C^k[a,b]$ 表示在闭区间上具有连续至 $k$ 阶导数的函数；端点导数按单侧延拓理解。$\|f\|_\infty=\max_{x\in[a,b]}|f(x)|$。向量和矩阵范数在第~\ref{sec:na-norms} 节定义。

$x_*$ 或 $x^*$ 表示精确解，$x_k$ 表示第 $k$ 次迭代，$\widehat x$ 表示最终计算值，$e_k=x_k-x_*$ 表示有符号误差。网格宽度记为 $h$，时间步长记为 $\Delta t$ 或在明确的时间章节中记为 $h$，单位舍入误差记为 $u$。$E(h)=O(h^p)$ 表示存在与充分小的 $h>0$ 无关的常数 $C$，使 $|E(h)|\le Ch^p$；常数可依赖固定区间、函数或问题参数，若这些参数也变化，需另行说明。

算法描述在未注明时采用精确算术；实际程序需增加有限值检查、迭代上限和失败状态。给定容差不是给定真误差：只有附有误差界或满足相应条件的估计量，停止判据才具有误差含义。计算代价中的 $O(n^p)$ 以浮点加、减、乘、除为基本操作计数，忽略常数；函数评价昂贵时应另计其次数。

